% ====================================================================
\section*{2.\quad Curvature and the genus tower}
% ====================================================================

\subsection*{The genus-$g$ propagator}

At genus $g\ge 1$, $z_1-z_2$ has no global meaning: the
curve~$X$ is compact, and any meromorphic function with a single
simple zero must also have a simple pole.

The replacement is the \emph{prime form} $E(z_1,z_2)$: a section
of $K^{-1/2}\boxtimes K^{-1/2}$ on $X\times X$, vanishing to
first order along the diagonal $\Delta$, with no other zeros.
The bundle is $K^{-1/2}\boxtimes K^{-1/2}$, \emph{not}
$K^{+1/2}\boxtimes K^{+1/2}$: the prime form is a
$(-\frac12,-\frac12)$-differential. In a local coordinate
where the diagonal is $z_1=z_2$,
\[
E(z_1,z_2)
\;=\;
(z_1-z_2)\cdot\bigl(1+O(z_1-z_2)\bigr)
\cdot(dz_1)^{-1/2}(dz_2)^{-1/2}.
\]

The prime form has nontrivial monodromy: analytic continuation
around the $B$-cycle~$B_j$ in~$z_1$ multiplies~$E$ by
$-\exp(-\pi i\Omega_{jj}-2\pi i\int_{z_2}^{z_1}\omega_j)$,
where $\Omega=(\Omega_{jk})$ is the period matrix. The
$A$-cycle monodromy is trivial. The \emph{genus-$g$ propagator}
compensates the $B$-cycle monodromy by a period correction:
\[
\eta^{(g)}_{12}
\;=\;
d\log E(z_1,z_2)
\;+\;
\pi\sum_{j,k=1}^{g}
(\mathrm{Im}\,\Omega)^{-1}_{jk}\,
\mathrm{Im}\!\Bigl(\int_{z_2}^{z_1}\omega_j\Bigr)\,
\omega_k(z_1).
\]
The correction renders $\eta^{(g)}_{12}$
single-valued on~$X$ at the cost of holomorphicity: the
correction term is real-analytic. The propagator couples
to the Hodge bundle
$\mathbb E\to\overline{\mathcal M}_g$ through the period
matrix~$\Omega$.

\subsection*{Curvature}

At genus~$0$, the Arnold relation
$\eta_{12}\wedge\eta_{23}
+\eta_{23}\wedge\eta_{31}
+\eta_{31}\wedge\eta_{12}=0$
holds identically. At genus $g\ge 1$, the period corrections
produce a residual $(1,1)$-form: the fibrewise bar differential satisfies
\begin{equation}\label{eq:pref-curvature}
d_{\mathrm{fib}}^{\,2}\;=\;\kappa(\cA)\cdot\omega_g,
\end{equation}
where $\omega_g$ is the Arakelov $(1,1)$-form on the fibre of the
universal curve $\pi\colon\mathcal C_g\to\overline{\mathcal M}_g$,
and $\kappa(\cA)\in\Bbbk$ is the scalar \emph{modular
characteristic}, determined entirely by the leading OPE singularity.

For the three principal families:
\begin{itemize}
\item \emph{Heisenberg} at level~$k$: the only OPE datum is
$\alpha_{(1)}\alpha=k$; every residue at a collision
stratum extracts this coefficient.
$\kappa(\mathcal H_k)=k$.

\item \emph{Kac--Moody}: the bracket
$J^a_{(0)}J^b=f^{ab}{}_{c}J^c$ and the inner product
$J^a_{(1)}J^b=k\kappa^{ab}$ both contribute. The
curvature receives the adjoint trace:
\[
\kappa(\widehat{\mathfrak g}_k)
\;=\;
\frac{(k+h^\vee)\dim\mathfrak g}{2h^\vee}\,.
\]
At the critical level $k=-h^\vee$, $\kappa=0$: the bar complex
is uncurved.

\item \emph{Virasoro} at central charge~$c$:
the bar propagator $d\log E(z,w)$ absorbs one pole order
(the $d\log$ measure sends $z^{-n}$ to $z^{-(n-1)}$), so the
collision residue $r(z)=(c/2)/z^3+2T/z$ has pole orders one
less than the OPE.
$\kappa(\mathrm{Vir}_c)=c/2$.
The Koszul dual is $\mathrm{Vir}_{26-c}$, with
$\kappa+\kappa^!=13$ (not zero); at $c=13$
the algebra is self-dual.
\end{itemize}

\subsection*{Period integrals restore nilpotence}

The curvature~\eqref{eq:pref-curvature} obstructs
$d_{\mathrm{fib}}^2=0$, but not the total $D_g^2=0$.
The corrected total differential
\[
D_g\;=\;d_{\mathrm{fib}}+\nabla^{\mathrm{GM}}
\]
incorporates the Gauss--Manin connection on the Hodge bundle.
The fibrewise curvature $\kappa\cdot\omega_g$ is a class in
$H^{1,1}(\mathcal C_g)$; the period map identifies this class
with the $(0,1)$-curvature of $\nabla^{\mathrm{GM}}$ on
$R^1\pi_*\mathbb C$. The two curvatures cancel:
\[
D_g^{\,2}
\;=\;
\kappa\cdot\omega_g-\kappa\cdot\omega_g+0
\;=\;0.
\]

\subsection*{The genus tower}

For uniform-weight algebras, the obstruction to flatness at
genus~$g$ is the characteristic class
\[
\mathrm{obs}_g(\cA)
\;=\;
\kappa(\cA)\cdot\lambda_g
\;\in\;
H^*(\overline{\mathcal M}_g),
\qquad
\lambda_g\;=\;c_g(\mathbb E).
\]
At genus~$1$ this holds unconditionally for all families; for
multi-weight families at genus~$g \geq 2$ the factorization
remains open
\textup{(}see \S\textup{\ref{sec:modular-char-spectral-disc-intro}}\textup{)}.
Grothendieck--Riemann--Roch on the universal curve computes
the total obstruction via the Todd class of the relative
tangent bundle, yielding the $\hat A$-genus:
\[
\sum_{g\ge 1}\mathrm{obs}_g(\cA)\cdot x^{2g}
\;=\;
\kappa(\cA)\cdot\bigl(\hat A(ix)-1\bigr),
\qquad
\hat A(ix)=\frac{x/2}{\sin(x/2)}
=1+\frac{x^2}{24}+\frac{7x^4}{5760}+\cdots\,.
\]
All coefficients are positive. The first values:
\[
F_1(\cA)
\;=\;
\frac{\kappa(\cA)}{24}\,,
\qquad
F_2(\cA)
\;=\;
\frac{7\,\kappa(\cA)}{5760}\,,
\qquad
F_3(\cA)
\;=\;
\frac{31\,\kappa(\cA)}{967680}\,.
\]

\subsection*{Complementarity}

Verdier duality splits the center local system into complementary
Lagrangian summands:
\[
R\Gamma(\overline{\mathcal M}_g,\;\mathcal Z_\cA)
\;\simeq\;
\mathbf Q_g(\cA)\;\oplus\;\mathbf Q_g(\cA^!).
\]
The natural pairing
$\mathbf Q_g(\cA)\otimes\mathbf Q_g(\cA^!)
\to\Bbbk[-(3g-3)]$
is a $(-(3g-3))$-shifted symplectic form, and the two summands
are Lagrangian.
At the scalar level, $\kappa(\cA)+\kappa(\cA^!)=0$
for Kac--Moody and free-field algebras; for $\mathcal W$-algebras
the sum is a nonzero constant ($13$ for Virasoro,
$250/3$ for~$\mathcal{W}_3$).

The Feigin--Frenkel involution $k\leftrightarrow -k-2h^\vee$
is Koszul duality on configuration spaces: it interchanges
$\cA$ and~$\cA^!$, and the identity
$\kappa(\widehat{\mathfrak g}_k)+
\kappa(\widehat{\mathfrak g}_{-k-2h^\vee})=0$
is a direct consequence. At the critical level, the Koszul dual
is the Langlands-dual critical-level algebra
$\widehat{{}^L\mathfrak g}_{-{}^Lh^\vee}$.


% ====================================================================
\section*{3.\quad The universal Maurer--Cartan element}
% ====================================================================

The scalar $\kappa(\cA)$ is only the leading term. The bar complex
carries a universal Maurer--Cartan element
$\Theta_\cA\in\mathrm{MC}(\mathfrak g^{\mathrm{mod}}_\cA)$
whose arity-$r$ projection yields a shadow invariant~$S_r(\cA)$.
Its construction requires three inputs: a
cooperadic chain complex on logarithmic configuration spaces (the
geometric input), a homotopy chiral algebra on the curve (the
algebraic input), and a convolution $L_\infty$-structure mediating
between them (the homotopical input).

\subsection*{The modular operad and the Feynman transform}

A \emph{modular operad}~$\cO$
(Getzler--Kapranov~\cite{GeK98}) has operations indexed
by \emph{connected stable graphs}: graphs with genus labels
$g(v)\ge 0$ at each vertex satisfying the stability condition
$2g(v)-2+\mathrm{val}(v)>0$, with total genus
$g(\Gamma)=b_1(\Gamma)+\sum_v g(v)$.
Trees are the genus-$0$ case. Stable graphs parametrise
the boundary stratification of~$\overline{\cM}_{g,n}$:
vertices are irreducible components, edges are nodes,
half-edges are marked points.

The \emph{Feynman transform} replaces trees in the
operadic bar construction by stable graphs:
\begin{equation}\label{eq:preface-feynman-sum}
F\cO(g,n)
\;=\;
\bigoplus_{\Gamma\in\mathsf{Gr}^{\mathrm{st}}_{g,n}}
\frac{\hbar^{g(\Gamma)}}{|\operatorname{Aut}(\Gamma)|}
\;\bigotimes_{v\in V(\Gamma)}\cO\bigl(g(v),\mathrm{In}(v)\bigr),
\end{equation}
with edge-contraction differential. The identity $d^2=0$ is
codimension-$2$ boundary cancellation: each pair of edges can be
contracted in either order, and the two terms cancel by the
associativity of~$\cO$.

When $\cO=\Com$ (the commutative modular operad, with
$\Com(g,n)=\Bbbk$ for every stable $(g,n)$), the Feynman
transform becomes the \emph{stable graph complex}
\[
F\Com(g,n)
\;=\;
\bigoplus_{\Gamma\in\mathsf{Gr}^{\mathrm{st}}_{g,n}}
\frac{\hbar^{g(\Gamma)}}{|\operatorname{Aut}(\Gamma)|}
\;\det\bigl(E(\Gamma)\bigr)\!:
\]
a chain complex spanned by isomorphism classes of stable graphs,
graded by edge number.

\subsection*{The convolution algebra}

At each level of generality, a bar-type construction is the
coefficient of a convolution dg Lie algebra whose Maurer--Cartan
elements classify algebraic structures:
\begin{equation}\label{eq:preface-convolution-ladder}
\begin{array}{r@{\;:\quad}c@{\qquad}l}
\text{algebra} &
 \Hom\bigl(B(\Ass),\,\End_A\bigr)
 & A_\infty\text{-structures on }A \\[4pt]
\text{operad} &
 \Hom_\Sigma\bigl(B(\cP),\,\End_V\bigr)
 & \cP_\infty\text{-structures on }V \\[4pt]
\text{modular} &
 \Hom_\Sigma\bigl(F\Com,\,\End_{\barB(\cA)}\bigr)
 & \text{consistent genus towers}
\end{array}
\end{equation}
The third line governs this book. Given a chiral algebra~$\cA$
with bar complex~$\barB(\cA)$, the \emph{modular convolution
$L_\infty$-algebra}
\begin{equation}\label{eq:pref-convolution}
\mathfrak g^{\mathrm{mod}}_\cA
\;:=\;
\Convinf\!\Bigl(
\cC^{\log\mathrm{FM}}_{\mathrm{mod}},\;
\mathrm{End}_{\mathrm{Ch}_\infty}(A^{\mathrm{ch}}_\infty)
\Bigr)
\end{equation}
is assembled from the log-FM chain cooperad
$\cC^{\log\mathrm{FM}}_{\mathrm{mod}}$
(chains on Mok's logarithmic Fulton--MacPherson spaces, with
cocomposition from inclusion of boundary strata) and the
$\mathrm{Ch}_\infty$-endomorphism operad of the homotopy chiral
algebra (the \v{C}ech totalisation of~$\cA$ on an affine cover,
which carries a $\mathrm{Ch}_\infty$-structure via the secondary
Borcherds operations).

The construction has a \emph{strict chart}: the coderivation dg Lie
algebra $\Convstr(\cC^{\log\mathrm{FM}}_{\mathrm{mod}},
\mathrm{End}_{\mathrm{Ch}_\infty})$, with classical Lie bracket and
no higher $L_\infty$ brackets, but the same Maurer--Cartan moduli.
The strict chart suffices for constructing the universal MC element;
the full $L_\infty$ structure is needed for transfer, formality,
and gauge equivalence.

\subsection*{The total differential and $D^2=0$}

The modular bar functor at type $(g,n)$ sums over stable
graphs~$\Gamma$: at each vertex, the transferred chiral
bar contribution; along each internal edge, the propagator;
the geometric factor from the log-FM stratum; and the orientation
line $\det(E(\Gamma))$ recording signs. The total differential
decomposes into six components:
\[
D^{\log}_{\mathrm{mod}}
\;=\;
\underbrace{d_{\check C}}_{\text{\v Cech}}
\;+\;
\underbrace{d_{\mathrm{Ch}_\infty}}_{\text{transferred chiral}}
\;+\;
\underbrace{d_{\mathrm{coll}}}_{\text{collision}}
\;+\;
\underbrace{d_{\mathrm{sew}}}_{\text{separating}}
\;+\;
\underbrace{d_{\mathrm{pf}}}_{\text{planted-forest}}
\;+\;
\underbrace{\hbar\,\Delta}_{\text{genus-raising}}.
\]
The first two act on vertex labels; the collision differential
contracts internal edges by OPE residues; the last three
encode positive-genus boundary operators (separating clutching,
planted-forest corrections from Mok's log-FM degeneration, and
non-separating self-gluing that raises genus by~$1$).

The nilpotence $(D^{\log}_{\mathrm{mod}})^2=0$ follows from three
mechanisms: the $\mathrm{Ch}_\infty$-algebra axioms on vertex labels,
the Leibniz rule for coderivation/derivation pairs on one-edge terms,
and codimension-$2$ cancellation on two-edge terms.
The third is the nontrivial one.
The normal-crossings property of Mok's log-FM compactification
(Mok, Theorem~3.3.1) guarantees that every codimension-two
intersection is transverse: each pair of boundary faces composes
with opposite orientations, and the two contributions cancel.
This is $\partial^2=0$ for a manifold with corners, lifted
through the residue correspondence.


\subsection*{The bar-intrinsic construction}

The element $\Theta_\cA$ is extracted from a differential that
already squares to zero. No equation is solved.

The total bar differential
$D_\cA\colon\mathrm{B}^{\mathrm{mod}}(\cA)\to
\mathrm{B}^{\mathrm{mod}}(\cA)$ acts on the modular bar complex.
Decompose $D_\cA=d_0+\Theta_\cA$, where
$d_0=d_{\mathrm{int}}+[\tau,-]$ is the tree-level local part
(internal differential of~$\cA$ plus commutator with the
tautological twisting morphism). The \emph{universal
Maurer--Cartan element} is the remainder:
\begin{equation}\label{eq:pref-theta}
\Theta_\cA
\;:=\;
D_\cA-d_0
\;=\;
\sum_{\Gamma\in\mathsf{Gr}^{\mathrm{st}}}
\frac{W^{\log\mathrm{FM}}_\Gamma}{|\operatorname{Aut}(\Gamma)|}
\otimes\Phi^\cA_\Gamma
\;\in\;
\operatorname{MC}(\gAmod).
\end{equation}
Here $W^{\log\mathrm{FM}}_\Gamma\in
H^*(\overline{\cM}^{\log}_{g(\Gamma),n(\Gamma)})$ is the
fundamental class of the boundary stratum indexed by~$\Gamma$,
and $\Phi^\cA_\Gamma\in\operatorname{Hom}(\bar\cA^{\otimes
n(\Gamma)},\Bbbk)$ is the chiral amplitude
(propagators on internal edges, operations at vertices,
cyclic trace at the output).

The MC property is immediate: $D_\cA^{\,2}=0$ by the geometric
fact that $\overline{\cM}_{g,n}^{\log}$ has normal-crossings
boundary; $d_0^2=0$ by construction; expanding
$(d_0+\Theta_\cA)^2=0$ gives
\[
d_0\Theta_\cA+\Theta_\cA\,d_0
+\Theta_\cA^2=0
\qquad\Longleftrightarrow\qquad
[d_0,\Theta_\cA]+\tfrac12[\Theta_\cA,\Theta_\cA]=0.
\]
The construction requires no restriction on~$\cA$: no
simple-Lie hypothesis, no convergence condition, no choice of
basis. The MC property is a formal consequence of $\partial^2=0$
on the moduli space.

\subsection*{Five projections of $\Theta_\cA$}

The five main theorems collectively characterize~$\Theta_\cA$:
its existence, its completeness as an invariant of~$\cA$,
and its structural content.

\smallskip

\emph{Theorem~A} (Bar-cobar adjunction).
The Ran-space bar functor $\barB_X$ and the cobar
functor~$\Omega_X$ form an adjunction
$\barB_X\dashv\Omega_X$ on the category of factorization
algebras and coalgebras on~$X$.
The Verdier intertwining $D_{\mathrm{Ran}}(\barB(\cA))\simeq
\barB(\cA^!)$ identifies the Verdier dual of the bar coalgebra
with the bar of the Koszul dual.
This is the arena in which~$\Theta_\cA$ lives.

\smallskip

\emph{Theorem~B} (Bar-cobar inversion).
On the Koszul locus, $\Omega(\barB(\cA))\xrightarrow{\sim}\cA$:
the cobar of the bar recovers the original algebra.
$\Theta_\cA$ is a complete invariant.

\smallskip

\emph{Theorem~C} (Complementarity).
The center local system decomposes
$R\Gamma(\overline{\mathcal M}_g,\mathcal Z_\cA)\simeq
\mathbf Q_g(\cA)\oplus\mathbf Q_g(\cA^!)$ into Lagrangian summands
under a $(-(3g-3))$-shifted symplectic form.
Verdier duality decomposes~$\Theta_\cA$ into complementary halves.

\smallskip

\emph{Theorem~D} (Modular characteristic).
For uniform-weight modular Koszul algebras,
$\mathrm{obs}_g(\cA)=\kappa(\cA)\cdot\lambda_g$ at all genera,
with $\hat A$-genus generating function. The leading coefficient
of~$\Theta_\cA$ is~$\kappa(\cA)$: universal, additive, and
duality-constrained.

\smallskip

\emph{Theorem~H} (Chiral Hochschild).
The chiral Hochschild cohomology $\mathrm{ChirHoch}^*(\cA)$
is polynomial in degrees $\{0,1,2\}$ and Koszul-functorial.
The coefficient ring of~$\Theta_\cA$ is the deformation ring
of~$\cA$.

\medskip

These five theorems are five projections of a single object. The
existence of~$\Theta_\cA$ is encoded by the bar-cobar adjunction
(Theorem~A); inversion proves it is a complete invariant
(Theorem~B); Verdier duality decomposes it into complementary halves
(Theorem~C); the leading coefficient extracts~$\kappa$ with
$\hat A$-genus generating function (Theorem~D); the coefficient ring
is the chiral Hochschild complex (Theorem~H).

\subsection*{The shadow obstruction tower}

The finite-order projections of~$\Theta_\cA$ form a Postnikov tower
$\Theta_\cA^{\leq 2} \to
\Theta_\cA^{\leq 3} \to
\Theta_\cA^{\leq 4} \to \cdots$
whose inverse limit
$\Theta_\cA = \varprojlim_r \Theta_\cA^{\leq r}$
converges unconditionally (no analytic input, guaranteed by
the completeness of the depth filtration).
On each primary deformation line~$L$, the shadow generating
function $H(t)=\sum_{r\ge 2}rS_r\,t^r$ satisfies the algebraic
relation
\begin{equation}\label{eq:preface-riccati-hook}
H(t)^2 \;=\; t^4\,Q_L(t),
\qquad
Q_L(t) = 4\kappa^2+12\kappa\alpha\,t+(9\alpha^2+16\kappa S_4)\,t^2,
\end{equation}
a quadratic polynomial in three genus-$0$ invariants
$(\kappa,\alpha,S_4)$:
the infinite tower is algebraic of degree~$2$. The discriminant
$\Delta = 8\kappa S_4$ governs the qualitative behaviour:
$\Delta=0$ (tower terminates, the Gaussian archetype) versus
$\Delta\neq 0$ (infinite but algebraic, the nonlinear regime).

\bigskip

\noindent
The modular characteristic $\kappa(\cA)$ is only the leading term.
The higher-arity projections of~$\Theta_\cA$ see the algebraic
complexity invisible to~$\kappa$: the Heisenberg algebra and the
Virasoro algebra have different~$\kappa$ but also different internal
structure: one terminates at tensor degree~$2$, the other generates
an infinite tower. These projections are governed by a single
algebraic relation.
